Give the definition of a tensor algebra

Provide the necessary constructions of the tensor algebra

Naturally the tensor product built on \(M \times N\) can be extended to \(\bigtimes_{i = 1}^{p} M_i\),
which we denote in the case that \(M_i = M_j\) for all \(i, j\) as \(\otimes^{p} M\).
Again for each module \(N\) and \(p\)-linear/multilinear map \(\phi : \bigtimes_{i = 1}^{p} M_i \to N \) there is a unique \(\Phi : \otimes^{p} M \to N\)
through which \(\phi\) factors \(\phi = \Phi \circ \pi\), where \(\pi\) is the canonical projection into the tensor product.
The special cases are 

\(\otimes^{0} M = A\), where \(A\) is the underlying coefficient ring.
\(\otimes^{1} M = M\)

\(x_1 \otimes x_2 \dots \otimes x_p = \pi(x_1, x_2, \dots, x_p) \) and elements of this shape
are called \(p\)-th order tensors. If \(x \in \otimes^{p} M\) can be written as a single elementary
/\(p\)-th order tensor then it is called decomposable. Consider \(\mathbb{R}^{2}\) then \((e_1 \otimes e_2) + (e_2 \otimes e_1)\)
is not decomposable (assume \((e_1 \otimes e_2) + (e_2 \otimes e_1) = (a_1 e_1 + a_2 e_2) \otimes (b_1 e_1 + b_2 e_2)\) after applying
distributivity there will be contradicting assertions for the products of \(a_ib_j\)), but \(e_1 \otimes e_2\) is.

In particular if \(M\) is free with basis \(\{m_i\}_{i \in I}\) then \(\{m_{i_1} \otimes m_{i_2} \dots m_{i_p}\}_{i_j \in I}\)
is a basis for \(\otimes^{p} M\) which is free of rank \(|I|^p\).

One can then define a multiplication for a \(p\)-th order x and \(q\)-th order y in \(\otimes^{p + q} M\) as \(x \otimes y = x_1 \otimes \dots x_p \otimes y_1 \otimes \dots y_1\) and more generally on \(\bigotimes M = \bigoplus_{p \geq 0} \otimes^{p} M\) as \(x y = \sum_{p \geq 0} \sum_{r + s = p} x_r \otimes y_s\), where \(x_r\), \(y_s\) are the \(r\)-th, \(s\)-th order terms in \(x, y\) (which are a sum of these elementary tensors).


Define the exterior algebra and prove its existence

A \(p\)-th linear map \(f : \bigtimes^{p} M \to N\) is called alternating if \(f(x_1, \dots, x_n) = 0\), whenever \(x_i = x_j\) for some \(i \neq j\).
In particular if \(2\) is not a zero divisor in the underlying ring \(A\), then \(f\) is alternating iff. 
\(f(x_{\sigma(1)}, \dots, x_{\sigma(p)}) = \text{sgn}(\sigma)f(x_1, \dots, x_p) = 0\) for all \(\sigma \in S_p\).

The exterior algebra \(\bigwedge^{p} M\) satisfies for each alternating \(f : \bigtimes^{p} M \to N\), there is some \(F : \bigwedge^{p} M \to N\)
through which \(f\) factors, \(f = F \circ \pi\).

Its existence can be guaranteed since \(N^p(M) \subset \otimes^{p} M\) with \(N^p(M)\) being generated by \(\{ x_1 \otimes \dots x_p | \exists x_i = x_j, (i \neq j) \}\) 
is indeed a submodule and \(\bigwedge^{p} M = \otimes^{p} M / N^p(M)\) and more generally \(\bigwedge M = \otimes M / N(M)\), since \(N(M)\) is an ideal
in the (infinite) tensor algebra.

If \(M\) is a free module with basis \(\{m_i\}_{i \in I}\), then \(\bigwedge^{p} M = 0\) for \(p > n\). For \(1 \leq p \leq n\) 
\(m_{i_1} \wedge \dots \wedge m_{i_p}\), \(1 \leq i_1 < \dots < i_p \leq n\) is a basis of \(\bigwedge^{p} M\) which is of rank \(\frac{n!}{p!(n - p)!}\)


Define the symmetric algebra and prove its existence

A \(p\)-th linear map \(f : \bigtimes^{p} M \to N\) is called symmetric if \(f(x_1, \dots, x_n) = f(x_{\sigma(1)}, \dots, x_{\sigma(p)})\), for all \(\sigma \in S_p\).

Similarly to the exterior product the symmetric product \(S^p(M) = \text{Sym}^{p}(M)\) 
satisfies for each symmetric \(f : \bigtimes^{p} M \to N\), there is some \(F : \text{Sym}^{p} M \to N\)
through which \(f\) factors, \(f = F \circ \pi\).

Its existence can be guaranteed since \(L^p(M) \subset \otimes^{p} M\) with \(L^p(M)\) generated by \(\{ x_1 \otimes \dots x_p - x_{\sigma(1)}, \dots, x_{\sigma(p)} | \sigma \in S_p \}\), 
is indeed a submodule and \(\text{Sym}^{p}(M) = \otimes^{p} M / L^p(M)\) and more generally \(\text{Sym}(M) = \otimes M / L(M) = \otimes / \bigoplus_{p \geq} L^{p} \), since \(L(M)\) is a (two-sided) ideal.

If \(M\) is a free module with basis \(\{m_i\}_{i \in I}\), then
\(m_{i_1} \cdot \dots \cdot m_{i_p}\), \(1 \leq i_1 < \dots < i_p \leq n\) is a basis of \(\bigwedge^{p} M\), where \(\cdot\) is the notation for the residual class of \(x_1 \otimes x_2 \dots\), and the module is of rank \(\frac{(n + p - 1)!}{(n-1)!(p)!}\) ((n + p - 1) over (n -1))

Further for a free module \(M\) of rank \(n\) there is a natural isomorphism between \(\text{Sym}(M) \simeq A[x_1, \dots, x_n]\),
namely \(\phi(x_1 \cdot x_2 \cdot \dots x_p) = x_1 x_2 \dots x_p\) (note that \(\text{Sym}(M)^{0} = A\)) and multiplication of homogenous
polynomials gives the equivalent multiplication of the residual classes of the tensors.


Proof the fundamental theorem of Galois theory

The following prerequisites are required:

Let \(\{\phi_i\}_{i \in I}\) be a set of distinct homomorphisms \(\phi_i : K \to K'\) both fields
and \(L = \{ a \in K | a = \phi_1(a) = \phi_2(a) \dots \}\) then \(L\) is a subfield of \(K\) and \([K : L] \geq |I|\).
The second is derived by creating a lin. sys. of eqs. of a basis for \(K\) as an \(L\)-v.s.
(each row contains only an \(|I|\) length sum of one basis element) and applying the \(\phi_i\) in each row
can be shown to give lin. dep. homomorphisms, but that's impossible due to ind. of characters (they're characters
from the multiplicative group \(K^{*}\) to \(K'^{*}\)).

It holds for finite subgroups \(H\) of \(\text{Aut}(K)\), \([K:\text{Fix}(K;H)] = |H|\) and \(\text{Aut}(K;\text{Fix}(K;H)) = H\) 
the latter following from the first, and the first is derived by a similar argument as for \([K : L] \geq |I|\) but one uses additionally 
the trace (the map \(a \mapsto \sum_{i \in I} \phi_i(a)\), non-zero for independent \(\phi_i\)) to show the required result.

Finally one needs to prove
\(\text{Fix}(K;\text{Aut}(K;L)) = L\) for all \(L \subset K\) subfields (say the \(K\) is a Galois extension of \(F \subset L \subset K\)).
Let \(G = \text{Aut}(K;F)\) and \(L' = \text{Fix}(K;\text{Aut}(K;L))\) then clearly \(L \subset L'\) to show the converse let \(\{\phi_j\}_{j \in J}\)
be a set of residual classes in \(G/H\) then \(\psi_i = \phi_i |_{L}\) are pairwise distinct as otherwise \(\phi_i, \phi_j\) would be the same equivalence class.
It holds \(L = \{ a \in L | a = \psi_1(a) = \psi_2(a) \dots \} = K\) with some effort, hence \([L : F] \geq |J|\) but
\(|G| = [K : L'][L' : L][L : F]\)
\(|G| = |G : H||H| = |J|[K : L']\) so it must be \([L' : L] = 1\).


Describe cyclotomic extensions

If \(p\) is a prime, the splitting field of \(X^p − 1\) over \(\mathbb{Q}\) is \(\mathbb{Q}[\zeta]\), where \(\eta\) is the root \(e^{\frac{2\pi i}{p}}\) . The
minimal polynomial of \(\zeta\) is \(X^{p−1} + X^{p−2} + \dots + X + 1\). So the degree of \(\mathbb{Q}[\zeta]\) over \(\mathbb{Q}\) is \(p − 1\)

The Galois group of \(\mathbb{Q}(\zeta_n)\) over \(\mathbb{Q}\) is isomorphic to the multiplicative group of integers modulo \( n \), \((\mathbb{Z}/n\mathbb{Z})^*\) of size \(n-1\):
\[ \text{Gal}(\mathbb{Q}(\zeta_n)/\mathbb{Q}) \cong (\mathbb{Z}/n\mathbb{Z})^* \]
An automorphism \(\sigma\) in this Galois group is determined by:
\[ \sigma(\zeta_n) = \zeta_n^k \]
for some \( k \in (\mathbb{Z}/n\mathbb{Z})^* \).

The degree of the cyclotomic field \(\mathbb{Q}(\zeta_n)\) over \(\mathbb{Q}\) is:
\[ [\mathbb{Q}(\zeta_n) : \mathbb{Q}] = \varphi(n) \]
where \(\varphi\) is the Euler totient function. (note that \(n\) is not prime)


This can be seen by
\(X^{n} - 1 = \Pi_{\zeta \in \Delta_n} (X - \zeta) = \Pi_{d|n}(\Pi_{\zeta \in \Lambda_d} (X - \zeta))\), where \(\Lambda_d\) is the set of \(d\)-th primitive roots of unity,
which are the elements in \(\Deta_n\), whose order is \(d\). We call

\(\Phi_n = \Pi_{\zeta \in \Lambda_n} (X - \zeta)\) the \(n\)-th cyclotomic polynomial and one can obtain a recursive formula
\(X^{n} - 1 = (\Pi_{d|n, d<n} \Phi_d)\Phi_n\) so \(\Phi_n = \frac{X^{n} - 1}{\Pi_{d|n, d<n} \Phi_d}\), starting with \(\phi_1 = X - 1, \phi_2 = X + 1\).
Since \(|\Lambda_n| = \varphi(n)\) (primitive roots!) and \(\Phi_n\) is the minimal polynomial of \(\zeta_n\), the result follows.

They have special significance due to Lagrange's resolvent, namely if \(K\) is a cyclic extension of \(F\) of degree \(d\)
and \(X^{d} - 1\) splits over \(F\), then \(K = F[\alpha]\) with \(\alpha^{d} \in F\).
The proof included that \(l = 1 + \zeta^{d-1}\sigma + \zeta^{d-2}\sigma + \dots \zeta^2 \sigma^{d - 2} + \zeta\sigma^{d-1}\) is not the zero map for some \(\beta \in K\)
and \(\alpha = l(\beta)\) is the primitive element in question, by determining that the minimal polynomial of \(\alpha\) has degree \(d\) and \(\sigma(\alpha)^{d} = \zeta^{d}\alpha^{d} = \alpha^{d}\). 


Define a radicals, radical tower and state Galois great theorem

A tower of fields
\(F = F_0 \subseteq F_1 \subseteq F_2 \subseteq \dots \subseteq F_k = K\)
is called a radical tower provided that for every index \(j\) there exists an \(\alpha_j\) in \(F_{j+1}\) and a positive
exponent \(n_j\) such that
\(F_{j+1} = F_j[\alpha]\) and \(\alpha_j \in F_j\)  .
In other words, each field \(F_{j+1}\) is an extension of its preceding field \(F_j\) by a single radical. The
top field \(K\) in a radical tower is called a radical extension of the bottom field \(F\).

If a polynomial f in \(F[X]\) splits over a radical extension of a field \(F\), we say
that \(f\) is solvable by radicals. In particular the roots of \(f\) may successively obtained
from the bottom field \(F\) by elementary arithmetic operations and taking successive roots of elements
in \(F\) (since if \(F_j[\alpha] = F_{j+1}\), there is some \(\alpha^{n_j} \in F_j\), which in turn can only
be already in \(F_{j-1}\) or \((\alpha^{n_j})^{n_{j-1}} \in F_{j-1}\), e.g. for quadratics \(F[\sqrt{b^2 - 4c}]\)

State and prove the existence of Galois extensions which are also radical.

Every radical extension \(K\) of \(F\) is a subfield of another extension \(L\) of \(F\) that
is both radical and Galois over \(F\).
\(L\) is the normal closure, as it is assumed \(\text{char} = 0\), there is a \(\gamma_1 \in K\), s.t. \(F[\gamma_1] = K\) 
(primitive element theorem). Let \(g\) be the minimal polynomial of \(\gamma_1\) and \(\gamma_i\) the remaining roots
one then constructs the tower
\(F \subset K = F[\gamma_1] \subseteq F[\gamma_1, \gamma_2] \subseteq \dots \subseteq L = F[\gamma_1, \gamma_2, \dots, \gamma_n]\)
This is not necessarily radical, but can be made so by choosing \(\{\alpha_k\} \subset K\) s.t. \(F \subseteq F[\alpha_1] \subseteq F[\alpha_2] \dots K\)
is radical and it can be shown that \(F[\gamma_1, \gamma_2, \dots, \gamma_{n-1}, \sigma(\alpha_1), \sigma(\alpha_2), \dots]\) is radical and since 
\(F[\alpha_1,\alpha_2,\dots,\alpha_k] = K = F[\gamma_1]\), we must have \(F[\sigma(\alpha_1), \sigma(\alpha_2),\dots,\sigma(\alpha_k)] = F[\sigma(\gamma_1)] = F[\gamma_n]\)
by transitivity of the \(\sigma \in \text{Gal}(K/F)\) and choosing a suitable \(\sigma\) from the start. Thus we have our radical tower.

For every radical extension \(K\) of \(F\) there exists a radical tower
\(F = E_0 \subseteq E_1 \subseteq E_2 \subseteq \dots \subseteq E_n = L\)
starting with \(F\) and such that
\(K\) is a subfield of \(L\),
\(L\) is Galois over \(F\),
each field \(E_j\) is Galois over its preceding \(E_{j−1}\) , and
each Galois group \(\text{Gal}(Ej/Ej−1)\) is abelian.

Note that we can always make a splitting field \(K\) of some polynomial \(f\) to a splitting field
of \((X^{n} - 1)f\) by adjoining \(\zeta_n\), which is actually a radical extension. Since splitting fields
of \((X^{n} - 1)\) are cyclic abelian, and since

Let \(n\) be a positive integer and \(F\) a field over which the polynomial \(X^n − 1\)
splits. If \(K\) is an extension of \(F\) then K is a cyclic extension
of \(F\), and the degree of \([K : F]\) divides \(n\).

we could instead consider the adjoined tower \(F \subseteq F[\zeta] \subseteq F_1[\zeta] \subseteq \dots K[\zeta]\)
giving the intuition why the theorem should be true.
this inspires the main theorem

A polynomial is solvable by radicals, then its Galois group is a solvable group.

Idea of proof:
Let \(K\) be the splitting field of \( f(x) \), which is solvable by radicals, so there exists a chain of field extensions:
\(
F \subset E_1 \subset E_2 \subset \cdots \subset L
\)
such that \(L\) is Galois over \(F\), each field \(E_j\) is Galois over its preceding \(E_{j−1}\)
and \(\text{Gal}(Ej/Ej−1)\) is abelian.
Additionally by the fundamental correspondence this introduces subgroups of \(\text{Gal}(L/F)\)
\(
\text{Gal}(L/L) \subseteq \text{Gal}(L/E_{n-1}) \subseteq \text{Gal}(L/E_{n-2}) \dots \text{Gal}(L/F)
\)
Thus \(\text{Gal}(L/F)\) is solvable iff. \(\text{Gal}(L/E_{i}) \unlhd \text{Gal}(L/E_{i-1})\) and \(\text{Gal}(L/E_{i-1}/\text{Gal}(L/E_{i})\)
is abelian, since \(\text{Gal}(L/E_{i-1}/\text{Gal}(L/E_{i}) \simeq \text{Gal}(E_{i-1}/E_{i})\) and homomorphisms of solvable groups are solvable
(as are subgroups) this is the case. Moreover \(\text{Gal}(L/F)\text{Gal}(L/K) \simeq \text{Gal}(K/F)\) so \(\text{Gal}(K/F)\) is solvable too,
which was to be shown.

With some considerable effort the converse can be proven too.
